\documentclass[11pt,twocolumn]{article}
\usepackage{setspace}
\singlespacing
\usepackage[margin=0.7in]{geometry}
\usepackage{times}
\usepackage{booktabs}
\usepackage{amsmath}
\usepackage{multirow}
\usepackage{hyperref}
\usepackage{url}
\usepackage{enumitem}
\usepackage{amssymb}
\setlist{nosep,leftmargin=*}

\title{\vspace{-2em}Reimplementing DoRA: Weight-Decomposed Low-Rank Adaptation}
\author{Eric Do (ed547) Levi Kronenthaler (ldk67) Eli Furgeson (ehf38)\\ \url{https://github.com/ehfurgeson/dora-repro}}
\date{}
\begin{document}
\maketitle
\thispagestyle{empty}

\vspace{-1em}
\section{Introduction}
\vspace{-1em}
Parameter-efficient fine-tuning (PEFT) makes it feasible to adapt large
language models on commodity hardware, but a persistent accuracy gap
separates PEFT from full fine-tuning (FT). \textbf{DoRA: Weight-Decomposed
Low-Rank Adaptation}~\cite{liu2024dora}, by Liu, Wang, Yin, Molchanov,
Wang, Cheng, and Chen (ICML 2024), is a drop-in replacement for
LoRA~\cite{hu2022lora} that closes much of that gap with no inference
overhead.

The paper's core insight is that an FT update can be decomposed into a
\emph{magnitude} (per-output-channel scalar) and a \emph{direction}
(unit-norm vector), and that LoRA's update couples these in a way that
diverges from FT. DoRA reparameterizes each pretrained linear weight as
$W = m \cdot V/\|V\|_c$, learns the direction $V$ with a low-rank update,
and trains $m$ directly. Across NLU, commonsense reasoning, and visual
instruction tuning, the authors report that DoRA matches or beats LoRA,
often at half the trainable-parameter budget, and merges back into a
plain linear at inference time.
\vspace{-1em}
\section{Chosen Result}
\vspace{-1em}
We target the commonsense-reasoning portion of Table~1 of the paper:
zero-shot accuracy on eight tasks (BoolQ, PIQA, SIQA, HellaSwag,
WinoGrande, ARC-e, ARC-c, OBQA) for Llama-2-7B \cite{touvron2023llama2} and Llama-3-8B \cite{llama3}, comparing
LoRA and DoRA across a range of ranks. We focus on the paper's two
headline claims: \emph{(i)} DoRA at rank $r$ matches LoRA at rank $2r$,
and \emph{(ii)} DoRA's accuracy is more rank-robust than LoRA's. Both
claims are reproducible on a single Colab A100, and together they
constitute the most-cited evidence that weight-decomposition is doing
the work, rather than just adding capacity.

\vspace{-1em}
\section{Methodology}
\vspace{-1em}
\paragraph{Implementation.}
\texttt{DoRALayer} (\texttt{code/dora.py}) is an \texttt{nn.Module}
wrapping an \texttt{nn.Linear}, holding (i)~frozen base weights $W$,
(ii)~trainable LoRA factors $A\in\mathbb{R}^{r\times d_{\text{in}}}$
(Kaiming-init) and $B\in\mathbb{R}^{d_{\text{out}}\times r}$
(zero-init), and (iii)~a trainable magnitude vector $m$ initialized to
the row-wise $\ell_2$ norm of $W$. The forward pass is\\
$W' = m \cdot \frac{W + \tfrac{\alpha}{r}\,BA}{\bigl\|W + \tfrac{\alpha}{r}\,BA\bigr\|_{\text{row}} + \varepsilon},\quad y = W'x + b.$\\
At step~0, $B{=}0$ and $m{=}\|W\|$, so $W'{=}W$ exactly, the layer is
identity-initialized. A \texttt{merge\_and\_unload} routine bakes $W'$
back into the underlying linear at inference time, so DoRA carries
\emph{zero} latency overhead post-training. LoRA is implemented via the
HuggingFace PEFT library~\cite{peft} for an apples-to-apples baseline,
using the same target modules, rank, and $\alpha$.
\vspace{-1em}
\paragraph{BiDoRA extension.}
As a small extension beyond the main DoRA reproduction, we implemented a BiDoRA-inspired training mode~\cite{qin2024bidora}. BiDoRA argues that DoRA's simultaneous optimization of magnitude $m$ and direction $\Delta V$ can over couple the two components, so it decouples them with a bi-level objective. Direction is learned on a training split while magnitude is updated using validation feedback. Our version is a lightweight approximation of this idea, we simply alternate optimizer steps between direction-only updates $(A,B)$ and magnitude-only updates $(m)$, which lets us test whether decoupling the two DoRA components changes training behavior.
\vspace{-1em}
\paragraph{Targets and hyperparameters.}
Following the paper we apply LoRA/DoRA to the attention projections
\texttt{q\_proj} and \texttt{v\_proj} only, with $\alpha=2r$. We sweep
both methods at $r\in\{4,8,16,32,64\}$ on both Llama-2-7B and Llama-3-8B
(20~runs total). Training uses batch~4 with gradient accumulation~4
(effective~16), bf16, lr $2{\times}10^{-4}$ with cosine schedule and
50-step warmup, on the Commonsense-170K instruction
set~\cite{hu2023llmadapters}.
\vspace{-1em}
\paragraph{Deviations from the paper.} Two compute-driven simplifications:
(a)~we cap training at \textbf{1{,}000 optimizer steps} ($\sim$16K
examples; about $1/10$ of one epoch) rather than the paper's 3 epochs,
and (b)~we restrict adapter targets to $\{q,v\}$ rather than the paper's
$\{q,k,v,\text{up},\text{down}\}$. We also use a single learning rate
across all trainable parameters, whereas the paper uses a separate
(typically $2{\times}$) learning rate for $m$ versus $BA$.
\vspace{-1em}
\paragraph{Evaluation.}
After training we run \texttt{merge\_and\_unload} and evaluate the
resulting plain HuggingFace model zero-shot via
\texttt{lm-eval-harness}~\cite{eval-harness} on seven of the paper's
eight tasks (we omit SIQA due to a harness configuration mismatch). The
identical protocol is applied to the merged-LoRA baselines.


Table~\ref{tab:results} reports per-task accuracy and the unweighted
mean across the seven tasks for all 20 runs.

\begin{table*}[t]
\centering
\footnotesize
\setlength{\tabcolsep}{4pt}
\begin{tabular}{llcccccccc}
\toprule
Base model & Method (rank) & BoolQ & PIQA & HellaSwag & WinoGrande & ARC-e & ARC-c & OBQA & \textbf{Avg} \\
\midrule
\multirow{10}{*}{Llama-2-7B}
 & LoRA $r{=}4$   & 78.50 & 79.05 & 75.70 & 71.82 & 69.91 & 40.87 & 43.60 & \textbf{65.64} \\
 & LoRA $r{=}8$   & 79.82 & 79.49 & 76.15 & 71.74 & 73.53 & 42.66 & 43.00 & \textbf{66.63} \\
 & DoRA $r{=}4$   & 77.74 & 79.05 & 75.55 & 72.38 & 70.71 & 41.89 & 44.00 & \textbf{65.90} \\
 & DoRA $r{=}8$   & 79.94 & 79.11 & 76.40 & 72.22 & 71.72 & 42.41 & 43.60 & \textbf{66.49} \\
\midrule
\multirow{10}{*}{Llama-3-8B}
 & LoRA $r{=}4$   & 82.75 & 81.45 & 78.77 & 74.11 & 76.98 & 50.77 & 44.60 & \textbf{69.92} \\
 & LoRA $r{=}8$   & 82.60 & 81.45 & 78.63 & 75.22 & 78.83 & 52.47 & 43.60 & \textbf{70.40} \\
 & DoRA $r{=}4$   & 83.39 & 81.56 & 78.63 & 76.48 & 80.18 & 54.01 & 43.40 & \textbf{71.09} \\
 & DoRA $r{=}8$   & 83.52 & 82.10 & 78.57 & 76.24 & 79.04 & 53.33 & 44.40 & \textbf{71.03} \\
\bottomrule
\end{tabular}
\caption{Zero-shot accuracy (\%) per task for low ranks $r\in\{4,8\}$. PIQA, HellaSwag, ARC-e, ARC-c, and OBQA use \texttt{acc\_norm}; BoolQ and WinoGrande use \texttt{acc}, following the harness defaults.}
\label{tab:results}
\end{table*}
\section{Results \& Analysis}

\vspace{-1em}
\paragraph{Headline claim reproduced on Llama-3.}
At low rank, DoRA outperforms LoRA on Llama-3-8B by margins consistent
with the paper. At $r{=}4$, DoRA scores 71.09 vs LoRA's 69.92
(\textbf{$+1.17$ pts}); at $r{=}8$, 71.03 vs 70.40 (\textbf{$+0.63$}).
The gap narrows at $r{=}16$ and reverses at $r{=}32$/$r{=}64$, where LoRA
narrowly leads. This matches the paper's framing: the decomposition is
most useful when capacity is scarce; once rank is large enough, plain
LoRA can recover the same accuracy.
\vspace{-1em}
\paragraph{Half-rank claim reproduced on Llama-3.}
The paper's signature claim is that DoRA at rank $r$ matches LoRA at
rank $2r$. We see exactly this pattern: DoRA $r{=}4$ (71.09) \emph{beats}
LoRA $r{=}8$ (70.40) by 0.69; DoRA $r{=}8$ (71.03) is within 0.14 of
LoRA $r{=}16$ (71.17); DoRA $r{=}16$ (70.62) is within 0.33 of LoRA
$r{=}32$ (70.95). DoRA achieves the same accuracy with half the
trainable parameters at every comparison point.
\vspace{-1em}
\paragraph{Rank robustness is the cleanest signal.}
Across $r\in\{4,8,16,32,64\}$, DoRA's average accuracy spreads only
\textbf{0.47 pts} on Llama-3 and \textbf{0.96 pts} on Llama-2. LoRA
spreads \textbf{1.25 pts} and \textbf{1.69 pts} respectively. DoRA is
essentially flat across a $16{\times}$ range of ranks on Llama-3
(70.62--71.09); LoRA needs $r{\geq}16$ to reach plateau and degrades
sharply below it. This is the qualitative match to the paper's
rank-robustness analysis (Fig.~4 in the original).
\vspace{-1em}
\paragraph{Llama-2 shows a muted version of the same effect.}
On Llama-2-7B the curve is qualitatively similar but the gap is smaller
(DoRA $+0.26$ at $r{=}4$, $-0.14$ at $r{=}8$). DoRA peaks at $r{=}16$
(66.86) and dips to 66.27 at $r{=}32$ before recovering --- the dip is
driven by ARC-e (71.68 vs 72.98 at $r{=}16$) and is most likely
seed/fine-tuning noise, not a structural artifact. We attribute the
muted overall effect to undertraining: 1{,}000 steps is $\sim$10\% of
one epoch on Commonsense-170K, vs the paper's three full epochs. The
DoRA advantage appears to scale with both model size and training
duration, both of which favor the paper's regime over ours.
\vspace{-1em}
\paragraph{No instability or merge bugs.}
Across all 20 runs we observed no NaNs, divergence, or merge-time
numerical drift; merged-vs-unmerged forward passes agreed to bf16
tolerance, including at the smallest rank ($r{=}4$) where LoRA's
update is most aggressively scaled.
\vspace{-1em}
\section{Reflections}
\vspace{-1em}
\paragraph{Lessons learned.}\textbf{The paper's claims are real but rank-dependent.} The DoRA
  $>$ LoRA gap is not a uniform 1-point lift: it lives almost entirely
  at low rank ($r{\leq}8$) and on larger base models. This is
  consistent with the paper's framing. \textbf{Rank robustness is the most stable signal.} Even where
  per-rank gaps were noisy, the spread across ranks consistently
  favored DoRA on both models. This suggests rank robustness, rather
  than absolute accuracy at a given rank, is the cleanest empirical
  fingerprint of the decomposition. \textbf{Identity initialization is a useful sanity check.}
  $B{=}0$, $m{=}\|W\|$ must reproduce the base model exactly at step~0;
  we caught one $\varepsilon$-placement bug this way. \textbf{The mathematical core fits in $\sim$15 lines of PyTorch.}
  Almost all of \texttt{train.py} is logistics (target discovery,
  state-dict serialization, merge semantics, optional Hub upload).
\vspace{-1em}
\paragraph{What we would do next.} Train at least one config for the full 3 epochs to confirm the
  DoRA gap widens at scale and to test whether the Llama-2 muted-effect
  hypothesis (undertraining) is correct. Extend targets to $\{q,k,v,\text{up},\text{down}\}$ and split the
  learning rate so $m$ trains at a different rate than $BA$,
  the two methodological deviations remaining from the paper's recipe. 
\vspace{-1em}
\paragraph{Future directions.} The magnitude/direction split is general,
applying it to other PEFT families (VeRA, IA$^3$), to convolutional
weights, or to diffusion-model adaptation are natural extensions. The
paper itself shows the decomposition composes with quantization
and with VeRA\cite{liu2024dora}.
\newpage
\begin{thebibliography}{9}
\footnotesize
\bibitem{liu2024dora}
S.-Y.~Liu, C.-Y.~Wang, H.~Yin, P.~Molchanov, Y.-C.~F.~Wang, K.-T.~Cheng, and M.-H.~Chen.
DoRA: Weight-Decomposed Low-Rank Adaptation.
\emph{ICML}, 2024.

\bibitem{hu2022lora}
E.~J.~Hu, Y.~Shen, P.~Wallis, Z.~Allen-Zhu, Y.~Li, S.~Wang, L.~Wang, and W.~Chen.
LoRA: Low-Rank Adaptation of Large Language Models.
\emph{ICLR}, 2022.

\bibitem{touvron2023llama2}
H.~Touvron et al.
Llama 2: Open Foundation and Fine-Tuned Chat Models.
arXiv:2307.09288, 2023.

\bibitem{llama3}
AI@Meta.
The Llama 3 Herd of Models. 2024.

\bibitem{peft}
S.~Mangrulkar, S.~Gugger, L.~Debut, Y.~Belkada, S.~Paul, and B.~Bossan.
PEFT: State-of-the-art Parameter-Efficient Fine-Tuning.
\url{https://github.com/huggingface/peft}, 2022.

\bibitem{qin2024bidora}
P.~Qin, R.~Zhang, and P.~Xie.
BiDoRA: Bi-level Optimization-Based Weight-Decomposed Low-Rank Adaptation.
arXiv:2410.09758, 2024.

\bibitem{hu2023llmadapters}
Z.~Hu, L.~Wang, Y.~Lan, W.~Xu, E.-P.~Lim, L.~Bing, X.~Xu, S.~Poria, and R.~K.-W.~Lee.
LLM-Adapters: An Adapter Family for Parameter-Efficient Fine-Tuning of Large Language Models.
\emph{EMNLP}, 2023.

\bibitem{eval-harness}
L.~Gao et al.
A framework for few-shot language model evaluation.
\url{https://github.com/EleutherAI/lm-evaluation-harness}, 2023.

\end{thebibliography}

\end{document}
